% Standalone problem document, generated from the adjacent README.md.
% Compile with XeLaTeX. Mathematical content is maintained in Markdown.
\documentclass[11pt,a4paper]{article}
\usepackage[fontsize=10.5pt]{scrextend}
\usepackage[margin=22mm,headheight=15pt,headsep=9mm,footskip=11mm]{geometry}
\usepackage{fontspec}
\setmainfont{texgyrepagella}[Extension=.otf,UprightFont=*-regular,BoldFont=*-bold,ItalicFont=*-italic,BoldItalicFont=*-bolditalic]
\setsansfont{texgyreheros}[Extension=.otf,UprightFont=*-regular,BoldFont=*-bold,ItalicFont=*-italic,BoldItalicFont=*-bolditalic]
\setmonofont{texgyrecursor}[Extension=.otf,UprightFont=*-regular,BoldFont=*-bold,ItalicFont=*-italic,BoldItalicFont=*-bolditalic,Scale=MatchLowercase]
\usepackage{amsmath,amssymb}
\usepackage{unicode-math}
\setmathfont{texgyrepagella-math.otf}
\usepackage{xcolor,microtype,enumitem,needspace,fancyhdr}
\usepackage{bookmark}
\definecolor{ink}{HTML}{17344B}
\definecolor{accent}{HTML}{197D80}
\definecolor{muted}{HTML}{596773}
\hypersetup{colorlinks=true,linkcolor=accent,urlcolor=accent,pdftitle={FR-11 - The
minimum number of quadratic measurements for generalized phase
retrieval},pdfauthor={Open Problems in Numerical Linear Algebra}}
\urlstyle{same}
\setlength{\parindent}{0pt}
\setlength{\parskip}{5pt plus 1pt minus 1pt}
\linespread{1.04}
\setlength{\emergencystretch}{3em}
\widowpenalty=10000
\clubpenalty=10000
\displaywidowpenalty=10000
\allowdisplaybreaks[1]
\setlist{nosep,leftmargin=1.5em}
\providecommand{\tightlist}{\setlength{\itemsep}{0pt}\setlength{\parskip}{0pt}}
\providecommand{\pandocbounded}[1]{#1}
\pagestyle{fancy}
\fancyhf{}
\fancyhead[L]{\sffamily\footnotesize\color{muted}OPEN PROBLEMS IN NUMERICAL LINEAR ALGEBRA}
\fancyhead[R]{\sffamily\bfseries\footnotesize\color{accent}FR-11}
\fancyfoot[L]{\parbox[b]{0.9\textwidth}{\sffamily\fontsize{8}{10}\selectfont\color{muted}Editorial ratings; a bounded search cannot certify that a problem remains unsolved.\\Literature check: 2026-09-11}}
\fancyfoot[R]{\sffamily\footnotesize\color{muted}\thepage}
\renewcommand{\headrulewidth}{0.3pt}
\renewcommand{\headrule}{\hbox to\headwidth{\color{accent}\leaders\hrule height \headrulewidth\hfill}}
\makeatletter
\renewcommand{\section}{\Needspace{5\baselineskip}\@startsection{section}{1}{0pt}{12pt plus 2pt minus 2pt}{4pt}{\sffamily\bfseries\large\color{ink}}}
\renewcommand{\subsection}{\Needspace{4\baselineskip}\@startsection{subsection}{2}{0pt}{9pt plus 1pt minus 1pt}{3pt}{\sffamily\bfseries\normalsize\color{ink}}}
\makeatother
\setcounter{secnumdepth}{0}
\begin{document}
{\sffamily\small\color{muted}Frames and matrix designs\par}
\vspace{3pt}
{\raggedright\hyphenpenalty=10000\exhyphenpenalty=10000\sffamily\bfseries\fontsize{21}{24}\selectfont\color{ink}The
minimum number of quadratic measurements for generalized phase
retrieval\par}
\vspace{7pt}
{\sffamily\small\textbf{Difficulty:} challenging\quad\textbf{Importance:} interesting
to the community\par}
{\sffamily\small\bfseries\color{accent}Status: Partially resolved\par}
\vspace{4pt}
{\color{accent}\hrule height 0.8pt}
\vspace{6pt}
\textbf{Rating rationale:} Exact thresholds across all dimensions
require algebraic and topological analysis beyond known families; they
guide quadratic matrix-measurement design.

\subsection{Reviewed submission -
2026-09-11}\label{reviewed-submission---2026-09-11}

Matthew J. Colbrook, Department of Applied Mathematics and Theoretical
Physics, University of Cambridge.

The construction theorem and its proof give twelve explicit integral
symmetric \(7\times7\) matrices, with maximum absolute entry 144, that
recover every real signal including zero up to sign. Thus
\(m_{\mathbb R}(7)\le12\) in the unrestricted self-adjoint model. No
matching lower bound or all-dimension/all-field formula is established.
The manuscript also proves quantitative separation and exact-arithmetic
decoding; floating-point decoder tests are finite diagnostics, with a
separately documented scaling correction.

See the
\href{https://github.com/ajt60gaibb/OpenProblemsInNLA/blob/main/references/colbrook-frames-2026-09-11/manuscripts/twelve_measurements.pdf}{manuscript},
\href{https://github.com/ajt60gaibb/OpenProblemsInNLA/blob/main/references/colbrook-frames-2026-09-11/verification/reviews/FR-11-review.md}{independent
agent review}, and
\href{https://github.com/ajt60gaibb/OpenProblemsInNLA/blob/main/references/colbrook-frames-2026-09-11/README.md}{reproducible
submission record}. Independent agent review is not external human peer
review or formal proof-assistant certification. Original problem,
ratings and historical audits are retained below.

\subsection{Problem statement}\label{problem-statement}

Fix a field \(\mathbb F\in\{\mathbb R,\mathbb C\}\) and an integer
\(d\ge2\). For self-adjoint matrices
\(A_1,\ldots,A_m\in\mathbb F^{d\times d}\) define \[
\Phi_A(x)=(x^*A_1x,\ldots,x^*A_mx)\in\mathbb R^m.
\] Call the family phase retrievable if, for every
\(x,y\in\mathbb F^d\), \[
\Phi_A(x)=\Phi_A(y)
\quad\Longrightarrow\quad
x=cy\text{ for some }c\in\mathbb F\text{ with }|c|=1.
\] Let \(m_{\mathbb F}(d)\) be the smallest \(m\) for which such a
family exists. Determine \(m_{\mathbb F}(d)\) exactly for all \(d\ge2\)
and both fields. The matrices may be indefinite and have any rank.
Recovery must hold for every signal, including zero. In the real case
the ambiguity is only sign. The source's real and complex versions are
grouped as one problem here.

\subsection{Relevance and ratings}\label{relevance-and-ratings}

This is the sharp measurement-design problem for recovering a rank-one
Gram matrix from linear matrix measurements. Allowing arbitrary
self-adjoint measurement matrices distinguishes it from the rank-one
random-frame question FR-05. Exact thresholds in all dimensions involve
substantial algebraic and topological obstructions.

\newpage
\subsection{References}\label{references}

\begin{itemize}
\tightlist
\item
  Z. Xu, \emph{Signal recovery on algebraic varieties using linear
  samples}, Acta Math. Sinica (English Series) 42 (2026), 740--754
  (\href{https://doi.org/10.1007/s10114-026-5299-y}{journal});
  \href{https://arxiv.org/pdf/2506.17572v2}{arXiv:2506.17572v2}, §4,
  Problems 4.1 and 4.3; Theorems 4.3 and 4.8 give exact values in some
  dimensions.
\item
  Y. Wang and Z. Xu, \emph{Generalized phase retrieval: Measurement
  number, matrix recovery and beyond}, Appl. Comput. Harmon. Anal. 47
  (2019), 423--446 (\href{https://arxiv.org/abs/1605.08034}{primary
  manuscript};
  \href{https://doi.org/10.1016/j.acha.2017.09.003}{journal}).
\end{itemize}

\subsection{Status check --- 2026-09-10}\label{status-check-2026-09-10}

Xu explicitly retains both full threshold questions. Bounds and exact
dimensional families do not determine either entire function. Searches
for generalized phase retrieval, minimum measurement number, and
2025--2026 found no complete solution. Xu's August 2026
\href{https://arxiv.org/abs/2608.15003}{almost-everywhere result}
concerns generic-signal recovery with rank-one measurements, not the
all-signals requirement here. This distinction was checked in the new
paper's stated theorem.

\textbf{Independent audit:} Independently rechecked Xu Problems 4.1 and
4.3 and the definitions preceding them. Theorems 4.3 and 4.8 settle
dimensional families inside the displayed target, so the appropriate
label is Partially resolved. Later generalized-phase-retrieval searches
found no all-dimension threshold formula.
\end{document}
